% !PRINTED 233
\weberpage{233}
%<<CONT
von denen die erste linear, die andere von der $n$ten Ordnung ist,
so bestimmen, daß

\weq{7}{\begin{pmatrix} \alpha', & \beta' \\ \gamma', & \delta' \end{pmatrix}\begin{pmatrix} a, & 0 \\ c, & \partial \end{pmatrix} = \begin{pmatrix} 1, & 0 \\ 0, & n \end{pmatrix}\begin{pmatrix} \alpha, & \beta \\ \gamma, & \delta \end{pmatrix}.}

Denn schreiben wir die Relation (7) so:
\[
\begin{pmatrix} \delta', & -\beta' \\ -\gamma', & \alpha' \end{pmatrix}\begin{pmatrix} 1, & 0 \\ 0, & n \end{pmatrix} = \begin{pmatrix} a, & 0 \\ c, & \partial \end{pmatrix}\begin{pmatrix} \delta, & -\beta \\ -\gamma, & \alpha \end{pmatrix},
\]
so leitet man daraus einfach her:

\weq{8}{\begin{aligned}
\delta' &= a\delta, & \gamma' &= \partial\gamma - c\delta,\\
\partial\beta' &= \beta, & n\alpha' &= \partial\alpha - c\beta, & a\alpha' &= \alpha - c\beta'.
\end{aligned}}

Hierdurch ist zunächst $\partial$ bestimmt als der \gesperrt{größte gemeinschaftliche
Teiler} von $\beta$ und $n$ oder [wegen (4)] von $\mu'$ und $n$. Denn wäre
$\partial$ nicht der größte gemeinschaftliche Teiler von $n=\partial a$ und
$\beta=\partial\beta'$, so müßten $a$ und $\beta'$ einen gemeinsamen Teiler
% [sic] The source prints \alpha=\alpha a'+c\beta, contrary to (8).
haben und folglich auch $\beta$ und $\alpha=\alpha a'+c\beta$, was nicht sein
kann. Dadurch ist zugleich $a=n/\partial$ bestimmt, und $c$ muß der Kongruenz

\weq{9}{\partial\alpha - c\beta \equiv 0\ (\mathrm{mod}\ n)}

genügen, wodurch jedoch $c$ nicht absolut, sondern nur nach dem Modul $a$
bestimmt ist. Man kann also $c$ z.\,B. noch an die Bedingung knüpfen, daß es
durch irgend eine Potenz von 2, oder wenn $a$ nicht durch 3 teilbar ist,
durch irgend eine Potenz von 3 teilbar sein soll, was wir später bisweilen
tun werden.

Die drei Zahlen $a$, $\partial$, $c$ können keinen gemeinsamen Teiler haben,
denn ein solcher wäre [nach (8)] auch Teiler von $\alpha$ und $\beta$, was
unmöglich ist. Bezeichnen wir also den größten gemeinsamen Teiler von $a$
und $\partial$ mit $e$, so muß $c$ \gesperrt{relativ prim} zu $e$ sein.

Die Zahlen $a$, $\partial$, $c$, letztere modulo $a$, sind wegen (4) und (8)
durch die beiden Zahlen $\mu$, $\mu'$ völlig bestimmt und ändern sich nicht,
wenn $\mu$, $\mu'$ mit einem gemeinsamen, zu $n$ teilerfremden Faktor
multipliziert, d.\,h. durch ein anderes Paar derselben \gesperrt{Reihe} ersetzt werden.
Zerlegt man $n$ irgendwie in zwei Faktoren $a$, $\partial$, so kann man für
$c$ noch
\[
\frac{a}{e}\varphi(e)
\]
nach dem Modul $a$ verschiedene zu $e$ teilerfremde Werte annehmen. Jede
dieser Annahmen führt zu einem Zahlenpaar $\mu$, $\mu'$ durch die
Kongruenzen
%CONT>>
